%% LyX 2.3.6.2 created this file.  For more info, see http://www.lyx.org/.
%% Do not edit unless you really know what you are doing.
\documentclass[12pt,english]{report}
\usepackage{lmodern}
\renewcommand{\familydefault}{\rmdefault}
\usepackage[T1]{fontenc}
\usepackage[latin9]{inputenc}
\usepackage{geometry}
\geometry{verbose,tmargin=0.5in,bmargin=0.5in,lmargin=0.5in,rmargin=0.5in}
\usepackage{fancyhdr}
\pagestyle{fancy}
\setcounter{secnumdepth}{3}
\setcounter{tocdepth}{3}
\usepackage{color}
\usepackage{babel}
\usepackage{float}
\usepackage{booktabs}
\usepackage{amsmath}
\usepackage{amssymb}
\usepackage{setspace}
\usepackage{esint}
\usepackage{microtype}
\onehalfspacing
\usepackage[unicode=true,pdfusetitle,
 bookmarks=true,bookmarksnumbered=false,bookmarksopen=false,
 breaklinks=true,pdfborder={0 0 0},pdfborderstyle={},backref=false,colorlinks=false]
 {hyperref}
\hypersetup{
 citecolor=blue}

\makeatletter

%%%%%%%%%%%%%%%%%%%%%%%%%%%%%% LyX specific LaTeX commands.
%% Because html converters don't know tabularnewline
\providecommand{\tabularnewline}{\\}
\floatstyle{ruled}
\newfloat{algorithm}{tbp}{loa}[chapter]
\providecommand{\algorithmname}{Algorithm}
\floatname{algorithm}{\protect\algorithmname}

%%%%%%%%%%%%%%%%%%%%%%%%%%%%%% Textclass specific LaTeX commands.
\numberwithin{equation}{section}

%%%%%%%%%%%%%%%%%%%%%%%%%%%%%% User specified LaTeX commands.
\usepackage[dvipsnames]{xcolor}
\usepackage{tikz}
\usetikzlibrary{calc}
\usepackage{anyfontsize}







%~~~~~~~~~~~~~~~~~~~~~Fancy ToC~~~~~~~~~~~~~~~~~~~~~~
\usepackage{titletoc}
\usepackage{titlesec}

\definecolor{doc}{RGB}{0,60,110}


\contentsmargin{0cm}
\titlecontents{chapter}[0pc]
{\addvspace{30pt}%
\begin{tikzpicture}[remember picture, overlay]%
\draw[fill=RoyalBlue,draw=RoyalBlue, rounded corners] (-4,-.1) rectangle (-0.15,.5);%
\pgftext[left,x=-2.7cm,y=0.2cm]{\color{white}\Large \chaptertitlename\ \thecontentslabel};%
\end{tikzpicture}\color{RoyalBlue}\large\bfseries}%
{}
{}
{\hspace*{6pt}\titlerule\hspace*{6pt}\large\bfseries \thecontentspage
\begin{tikzpicture}[remember picture, overlay]
\draw[fill=doc!25,draw=RoyalBlue, rounded corners=0pt] (2pt,0) rectangle (6,0.1pt);
\end{tikzpicture}}%
\titlecontents{section}[2.4pc]
{\addvspace{1pt}}
{\contentslabel[\thecontentslabel]{2.4pc}}
{}
{\hfill\small \thecontentspage}
[]
\titlecontents{subsection}[4.8pc]
{\addvspace{1.0pt}}
{\contentslabel[\thecontentslabel]{2.4pc}}
{}
{\hfill\small\thecontentspage}
[]
\pagestyle{empty}
\makeatletter
\renewcommand{\tableofcontents}{%
\chapter*{%
\vspace*{-20\p@}%
\begin{tikzpicture}[remember picture, overlay]%
\pgftext[right,x=15cm,y=0.2cm]{\color{RoyalBlue}\Huge \contentsname};%
\draw[fill=RoyalBlue,draw=RoyalBlue, rounded corners=15pt] (13,-.75) rectangle (20,1);%
\clip (13,-.75) rectangle (20,1);
\pgftext[right,x=15cm,y=0.2cm]{\color{white}\Huge \contentsname};%
\end{tikzpicture}}%
\@starttoc{toc}}

%~~~~~~~~~~~~~~~~~Hyperref Settings~~~~~~~~~~~~~~~~~~~~~~

\usepackage{hyperref}
\hypersetup{%
    pdfborder = {0 0 0},
    colorlinks,
    citecolor=red,
    filecolor=green,
    linkcolor=RoyalBlue,
    urlcolor=cyan!50!black!90
}


%......Sectioning...%
\def\chpcolor{RoyalBlue}
\def\chpcolortxt{RoyalBlue}
\def\sectionfont{\LARGE}

\setcounter{secnumdepth}{2}


%Section:
\def\@sectionstrut{\vrule\@width\z@\@height12.5\p@}
\def\@makesectionhead#1{%
  {\par\vspace{20pt}%
   \parindent 0pt\raggedleft\sectionfont
   \colorbox{\chpcolor}{%
     \parbox[t]{90pt}{\color{white}\@sectionstrut\@depth4.5\p@\hfill
       \ifnum\c@secnumdepth>\z@\thesection\fi}%
   }%
   \begin{minipage}[t]{\dimexpr\textwidth-90pt-2\fboxsep\relax}
   \color{\chpcolortxt}\@sectionstrut\hspace{5pt}#1
   \end{minipage}\par
   \vspace{10pt}%
  }
}
\def\section{\@afterindentfalse\secdef\@section\@ssection}
\def\@section[#1]#2{%
  \ifnum\c@secnumdepth>\m@ne
    \refstepcounter{section}%
    \addcontentsline{toc}{section}{\protect\numberline{\thesection}#1}%
  \else
    \phantomsection
    \addcontentsline{toc}{section}{#1}%
  \fi
  \sectionmark{#1}%
  \if@twocolumn
    \@topnewpage[\@makesectionhead{#2}]%
  \else
    \@makesectionhead{#2}\@afterheading
  \fi
}
\def\@ssection#1{%
  \if@twocolumn
    \@topnewpage[\@makesectionhead{#1}]%
  \else
    \@makesectionhead{#1}\@afterheading
  \fi
}

\usepackage{algorithm}
\usepackage[noend]{algpseudocode}
\usepackage{microtype}



\makeatletter
\newenvironment{breakablealgorithm}
  {% \begin{breakablealgorithm}
   \begin{center}
     \refstepcounter{algorithm}% New algorithm
     \hrule height.8pt depth0pt \kern2pt% \@fs@pre for \@fs@ruled
     \renewcommand{\caption}[2][\relax]{% Make a new \caption
       {\raggedright\textbf{\ALG@name~\thealgorithm} ##2\par}%
       \ifx\relax##1\relax % #1 is \relax
         \addcontentsline{loa}{algorithm}{\protect\numberline{\thealgorithm}##2}%
       \else % #1 is not \relax
         \addcontentsline{loa}{algorithm}{\protect\numberline{\thealgorithm}##1}%
       \fi
       \kern2pt\hrule\kern2pt
     }
  }{% \end{breakablealgorithm}
     \kern2pt\hrule\relax% \@fs@post for \@fs@ruled
   \end{center}
  }
\makeatother

\makeatother

\begin{document}
\pagestyle{empty}
\begin{tikzpicture}[overlay,remember picture]
\fill[black!2] (current page.south west) rectangle (current page.north east);
\begin{scope}[transform canvas ={rotate around ={45:($(current page.north west)+(-.5,-6)$)}}]
\shade[rounded corners=18pt, left color=Dandelion, right color=Dandelion!40] ($(current page.north west)+(-.5,-6)$) rectangle ++(9,1.5);
\end{scope}
\begin{scope}[transform canvas ={rotate around ={45:($(current page.north west)+(.5,-10)$)}}]
\shade[rounded corners=18pt, left color=lightgray,right color=lightgray!60] ($(current page.north west)+(0.5,-10)$) rectangle ++(15,1.5);
\end{scope}
\begin{scope}[transform canvas ={rotate around ={45:($(current page.north west)+(0.5,-10)$)}}]
\shade[rounded corners=8pt, left color=lightgray] ($(current page.north west)+(1.5,-9.55)$) rectangle ++(7,.6);
\end{scope}
\begin{scope}[transform canvas ={rotate around ={45:($(current page.north)+(-1.5,-3)$)}}]
\shade[rounded corners=12pt, left color=orange!80, right color=orange!60] ($(current page.north)+(-1.5,-3)$) rectangle ++(9,0.8);
\end{scope}
\begin{scope}[transform canvas ={rotate around ={45:($(current page.north)+(-3,-8)$)}}]
\shade[rounded corners=28pt, left color=red!80, right color=red!80] ($(current page.north)+(-3,-8)$) rectangle ++(15,1.8);
\end{scope}
\begin{scope}[transform canvas ={rotate around ={45:($(current page.north west)+(4,-15.5)$)}}]
\shade[rounded corners=25pt, left color=RoyalBlue
,right color=Emerald] ($(current page.north west)+(4,-15.5)$) rectangle ++(30,1.8);
\end{scope}
\begin{scope}[transform canvas ={rotate around ={45:($(current page.north west)+(13,-10)$)}},]
\shade[rounded corners=22pt, left color=orange, right color=Dandelion] ($(current page.north west)+(13,-10)$) rectangle ++(15,1.5);
\end{scope}
\begin{scope}[transform canvas ={rotate around ={45:($(current page.north west)+(18,-8)$)}},]
\shade[rounded corners=8pt, left color=lightgray] ($(current page.north west)+(18,-8)$) rectangle ++(15,0.6);
\end{scope}
\begin{scope}[transform canvas ={rotate around ={45:($(current page.north west)+(19,-5.65)$)}},]
\shade[rounded corners=12pt, left color=black] ($(current page.north west)+(19,-5.65)$) rectangle ++(15,0.8);
\end{scope}
\begin{scope}[transform canvas ={rotate around ={45:($(current page.north west)+(20,-9)$)}}]
\shade[rounded corners=20pt, left color=OrangeRed, right color=red!80] ($(current page.north west)+(20,-9)$) rectangle ++(14,1.2);
\end{scope}
\draw[ultra thick,gray] ($(current page.center)+(5,2)$) -- ++(0,-3cm) node[midway,left=0.25cm,text width=5cm,align=right,black!75]{{\fontsize{12}{12} \selectfont \bf AN \\ INTRODUCTION TO  \\[5pt]  BAYESIAN ANALYSIS}} node[midway,right=0.25cm,text width=6cm,align=left,orange]{{\fontsize{30}{50} \selectfont MTH535A}};
\node at ($(current page.center)+(0,-3)$) {{\fontsize{17}{17} \selectfont \bf A Comparative study of Bayesian approach to case-control studies}};
\node at ($(current page.center)+(0,-4)$) {{\fontsize{17}{17} \selectfont \bf  with errors in covariables}};
\node[text width=20cm,align=center] at ($(current page.center)+(0,-6.5)$) {{\fontsize{12}{12} \selectfont \textcolor{RoyalBlue}{ \bf Soumya Paul, Sankhadeep Mitra, Sampriti Dutta}} \\[3pt] \emph{Under the guidance of} \\ \selectfont \textcolor{Emerald}{Prof. Arnab Hazra} \\  Department of Mathematics and Statistics, \\[3pt] Indian Institute of Technology Kanpur};
\node at ($(current page.center)+(0,-10.5)$) {{\includegraphics[width=4cm,height=4cm,keepaspectratio]{IIT_Kanpur_Logo}}};
\end{tikzpicture}

\newpage{}

\tableofcontents{}

\pagebreak{}
\begin{abstract}
We would like to develop a Bayesian methodology for the analysis of
case-control data with covariate uncertainty. Pretending that the
distribution of imprecisely measured covariates is discrete on a heuristically
chosen support set provides a method that is relatively easy to implement
and applicable to a variety of study designs. Further development
of the method highlights the interplay between retrospective and prospective
analyses. We illustrate this method with simulated data.�
\end{abstract}

\section{Introduction}

A \emph{case-control} study is a study of a medical condition/disease
by comparing patients with the condition/disease (\emph{\textcolor{Emerald}{"cases"}})
to otherwise similar patients (\emph{\textcolor{Emerald}{"control"}})
who do not have the condition/disease. It is commonly used to identify
possible contributing factors. Case-control studies are an important
and useful method for studying health outcomes, and many methods have
been developed to analyze case-control data of exposed subjects between
cases and controls.�

Our analysis seeks to examine the effects of covariates on disease
status using samples collected as a function of disease status. In
many cases, some covariates cannot be measured precisely. It is also
well known that an analysis that treats inaccurate measurements as
accurate can lead to distorted results. In this paper, we presented
a Bayesian methodology to account for covariate measurement errors
in case-control analysis.�

As of late, Bayesian strategies have been effectively applied to a
wide run of measurment error issues \emph{\cite{Carroll1999,Dellaportas1995BAYESIANAO,Mallick1996,1995,Richardson_1993}.}
The victory of these strategies stems from a direct conceptual and
computational approach to averaging over unobserved quantities such
as the genuine covariate values, within the posterior distribution.�

If the precise covariates are denoted by $X$, and disease status
by $D$ (where $D=0$ and $D=1$ corresponds to disease free and diseased
respectively), then a case\textendash control study involves sampling
\textquotedblleft controls\textquotedblright{} from $X\mid D=0$ and
\textquotedblleft cases\textquotedblright{} from $X\mid D=1$. Therefore,
an exact analysis requires a likelihood based on the conditional distribution
of $X$ given $D$, the so-called \textbf{\emph{retrospective model}}. 

However, the standard approach to analysis is to pretend that the
data were collected prospectively and use�likelihood�based on the
\textbf{\emph{prospective model}} $D|X$ that is usually assumed to
follow a logistic regression model. A justification for this is provided
by\textcolor{yellow}{{} }\emph{\cite{PRENTICE_1979}}. In the case of
measurement error, the imprecise measurement $W$ stands in for $X$. 

Here we used an approach that derives the posterior distribution directly
from the retrospective model. In fact, as an approximation of the
retrospective posterior distribution, we get the�prospective�posterior
distribution. This sheds new light on the interplay of retrospective
and prospective analyses.�

This method can be applied to different measurement error scenarios.
One scenario is�\textbf{\emph{validation design}} where both $W$
and $X$ are measured for some subjects (full data) and only $W$
for the rest (reduced data). If the exact form of the distribution
of $W\mid X,D$ contains unknown parameters, the full data will contain
information about these parameters. Therefore, both full and reduced
data contain information about the relationship between $X$ and $D$.�

Another scenario arises when the distribution of $W\mid X,D$ is completely
known from an external study and the case-control study consists only
of reduced data. This situation is called \textbf{\emph{external validation}}\emph{.}
For our project, we will limit ourselves to external validation.�

The approach of \emph{\cite{Gustafson2000}} contains the�normal�retrospective
model for $X\mid D$ and the�normal�measurement error model for $W\mid X,D$.
The assumption of normal distribution is limited in scope. The approach
described here is much more widely applicable.�

\emph{Section 0.2} describes the model and derives some variants of
the posterior and prior distributions of interest. \emph{Section 0.5}
test the method on simulated data. \emph{Sections 0.3} and \emph{0.4}
details the Markov Chain Monte Carlo (MCMC) scheme used.

\section{Methodology}

Suppose for a particular individual,
\begin{center}
\[
D=\begin{cases}
1 & \textrm{if the individual is diseased}\\
0 & \textrm{if the individual is disease-free}
\end{cases}\textrm{ \textrm{and }}\mathbf{\boldsymbol{X}}=\left(X_{1},X_{2},....,X_{p}\right)\;\textrm{: a vector of \ensuremath{p} covariates.}
\]
\par\end{center}

Generally, the scale at which measurement error is additive may not
be same as the scale at which the covariate is related with the disease
model for imprecisely measured covariates $\left(\boldsymbol{W}\right)$.�Therefore,
we assume that there exists a suitable transformation $s\left(.\right)$
on $\boldsymbol{X}$ such that the measurement error has an additive
effect on the components of $\boldsymbol{X}$.

Our analysis requires the likelihood based on the retrospective model.
The standard approach to analysis, is to pretend the data are sampled
\emph{prospectively} and use a likelihood based on $D|\boldsymbol{X}$,
which is typically assumed to follow a \emph{logistic regression model}.
Thus we specify the disease's \emph{prospective model}� as

\begin{equation}
\log\left(\frac{\mathbb{P}(D=1|\boldsymbol{X}=\boldsymbol{x})}{\mathbb{P}(D=0|\boldsymbol{X}=\boldsymbol{x})}\right)=\alpha^{*}+\beta^{T}s\left(\boldsymbol{x}\right)\label{eq:2.1}
\end{equation}

, where $s\left(\boldsymbol{x}\right)=\left(s_{1}\left(x_{1}\right),s_{2}\left(x_{2}\right),....,s_{p}\left(x_{p}\right)\right)^{T}$,
with $s_{i}:\mathbb{R}\rightarrow\mathbb{R}$ being a known function.

We are interested in knowing the value of the parameter set $\left\{ \alpha^{*},\beta_{1},\beta_{2},....,\beta_{p}\right\} $. 

The data are collected retrospectively, i.e., the \emph{controls}
comprise a sample of size $n_{1}$ drawn from the distribution of
$\boldsymbol{X}\mid D=0$, while the \emph{cases} comprise a sample
of size $n_{2}$ drawn from the distribution of $\boldsymbol{X}\mid D=1.$
For future use we define the total sample size as $n=n_{1}+n_{2}$,
and the sampling fractions as $r_{i}={\displaystyle \frac{n_{i}}{n};\;i=1,2}$

In terms of log-odds there is a fundamental link between \emph{retrospective}
and \emph{prospective} model:

\begin{eqnarray}
\log\left(\frac{f\left(\boldsymbol{x}\mid D=1\right)}{f\left(\boldsymbol{x}\mid D=0\right)}\right) & = & \log\left(\frac{\mathbb{P}\left(D=1|\boldsymbol{X}=\boldsymbol{x}\right)f\left(\boldsymbol{x}\right)\mathbb{P}\left(D=0\right)}{\mathbb{P}\left(D=0|\boldsymbol{X}=\boldsymbol{x}\right)f\left(\boldsymbol{x}\right)\mathbb{P}\left(D=1\right)}\right)\nonumber \\
 & = & \log\left(\frac{\mathbb{P}\left(D=1|\boldsymbol{X}=\boldsymbol{x}\right)}{\mathbb{P}\left(D=0|\boldsymbol{X}=\boldsymbol{x}\right)}\right)-\log\left(\frac{\mathbb{P}\left(D=1\right)}{\mathbb{P}\left(D=0\right)}\right)\nonumber \\
 & = & \alpha^{*}+\beta^{T}s\left(\boldsymbol{x}\right)-\log\left(\frac{\mathbb{P}\left(D=1\right)}{\mathbb{P}\left(D=0\right)}\right)\nonumber \\
 & = & \alpha+\beta^{T}s\left(\boldsymbol{x}\right)\textrm{ where }\alpha=\alpha^{*}-\log\left(\frac{\mathbb{P}\left(D=1\right)}{\mathbb{P}\left(D=0\right)}\right)\label{eq:2.2}
\end{eqnarray}

Using the retrospective data we cannot estimate $\alpha^{*}$ without
knowing the value of $\mathbb{P}\left(D=1\right)$.

The equation (\ref{eq:2.2}) can be rewritten as equation which is
parametrized according to $\left(\alpha,\beta,h\right)$
\[
f_{\boldsymbol{X}\mid D}\left(\boldsymbol{x}\mid d\right)=\exp\left(d\left(\alpha+\beta^{T}s\left(\boldsymbol{x}\right)\right)\right)h\left(\boldsymbol{x}\right);\;\textrm{where \ensuremath{d=0,1} and \ensuremath{h\left(x\right):} density of \ensuremath{X\mid D=0}}
\]

If $d=1$, then we get the density of $\boldsymbol{X}\mid D$, i.e.,
\begin{eqnarray}
f_{\boldsymbol{X}\mid D}\left(x\mid d=1\right) & = & \exp\left(\left(\alpha+\beta^{T}s\left(\boldsymbol{x}\right)\right)\right)h\left(\boldsymbol{x}\right)\label{eq:2.3}\\
 & \textrm{and}\nonumber \\
\intop_{-\infty}^{\infty}f_{\boldsymbol{X}\mid D}\left(\boldsymbol{x}\mid d=1\right)d\boldsymbol{x} & = & 1\nonumber \\
\implies\intop_{-\infty}^{\infty}\exp\left(\alpha+\beta^{T}s\left(\boldsymbol{x}\right)\right)h\left(\boldsymbol{x}\right)d\boldsymbol{x} & = & 1\nonumber \\
\implies E_{h}\left(\exp\left(\alpha+\beta^{T}s\left(\boldsymbol{x}\right)\right)\right) & = & 1\label{eq:2.4}
\end{eqnarray}

We have equation (\ref{eq:2.3}) with constraint (\ref{eq:2.4}) and
three unknown parameters $\left(\alpha,\beta,h\right)$. So, $\alpha$
can be written as a function of $\beta$ and $h$ i.e. $\alpha=\alpha\left(\beta,h\right)$.

Let's talk about measurement error.Suppose we have only one covariate
$\left(p=1\right)$ for which the presence of measurement error is
externally validated i.e. we have observed $W$ where the actual value
$X$ is unobserved for all subjects and the measurement error density,
$f(w|x,d)$ is completely specified. Instead of restricting ourselves
to a specific form of the density, it is useful for what follows to
define $\tau^{2}=Var(W|X,D)$ as the variance of imprecise measurement
given actual measurement. Here, we presume that $D$ has no effect
on the conditional variance. More broadly, measurement error is frequently
considered to be \emph{nondifferential}, in which case the conditional
distribution of $W|X,D$ is independent of $D$. 

Now, the joint distribution of $\left(W,X,D\right)$ can be presented
as
\[
f\left(w,x,d\right)=f\left(w\mid x,d\right)f\left(x\mid d\right)f\left(d\right);\textrm{ where \ensuremath{f\left(x\mid d\right)} depends on \ensuremath{h} and \ensuremath{\beta} }
\]

For a given prior $f\left(\beta,h\right)$ the posterior distribution
of $\left(X,\beta,h\right)$ is
\begin{eqnarray*}
f(x,\beta,h|w,d) & \propto & f(w,x,d|\beta,h)f(\beta,h)\\
 & = & \left(\prod_{i=1}^{n}f\left(w_{i}\mid x_{i},d_{i}\right)f\left(x_{i}\mid d_{i},\beta,h\right)f\left(d_{i}\mid\beta,h\right)\right)f\left(\beta,h\right)
\end{eqnarray*}

Usually $\beta$ is the parameter of interest and density $h$\textcolor{yellow}{{}
}is the nuisance parameter. As per \emph{\cite{BUONACCORSI1990}},
If we assume $h$ belongs to a simple parametric family, there is
a risk of misspecifying the model. On the other hand, the flexible
model of $h$ proposed by \emph{\cite{Muller_1997,Muller_1999}} failed
to completely eliminate erroneous specifications, leading to complex
model-tuning procedures. To find a compromise between simplicity and
flexibility, here we consider approximating $h$ by a discrete distribution.
�

Instead of modeling $h$ itself as a discrete distribution, it makes
sense to reparameterize from $\left(\beta,h\right)$ to $\left(\beta,g\right)$,
where the density of $X$ is given by 
\begin{equation}
g\left(x\right)=r_{1}h\left(x\right)+r_{2}\exp\left\{ \alpha\left(\beta,h\right)+\beta s\left(x\right)\right\} h\left(x\right)\label{eq:2.6}
\end{equation}

Suppose $\lambda$ parameterizes $g$ and for simplicity, assume $\beta$
and $\lambda$ are a priori independent. Then (\ref{eq:2.5}) reduces
to
\begin{eqnarray}
f(x,\beta,\lambda|w,d) & \propto & \left(\prod_{i=1}^{n}f\left(w_{i}\mid x_{i},d_{i}\right)f\left(x_{i}\mid d_{i},\beta,\lambda\right)f\left(d_{i}\mid\beta,\lambda\right)\right)f\left(\beta\right)f\left(\lambda\right)\label{eq:2.7}
\end{eqnarray}

To approximate $g$ with a discrete distribution, we need to specify
a grid of points to act as supports.�Let's say the grid-points are
\[
z\left[1\right]<z\left[2\right]<........<z\left[m\right]
\]
, and define $\theta=\left\{ \theta_{1},\theta_{2},....,\theta_{n}\right\} $
to match the sample drawn from distribution of $X$ to the grid via
\[
x_{i}=z\left[\theta_{i}\right]\;;\theta_{i}\in\left\{ 1,2,.....,m\right\} 
\]

Further, let $\lambda_{j}$ be the probability that $g$ assigns to
$z\left[j\right]$, and say that $\lambda$ has a uniform, or $Dirichlet(1,...,1)$,
prior on the $m$-dimensional probability simplex. Under this formulation,
(\ref{eq:2.7}) becomes
\begin{eqnarray}
f(\theta,\beta,\lambda|w,d) & \propto & \left(\prod_{i=1}^{n}f\left(w_{i}\mid z\left[\theta_{i}\right],d_{i}\right)f\left(z\left[\theta_{i}\right]\mid d_{i},\beta,\lambda\right)f\left(d_{i}\mid\beta,\lambda\right)\right)f\left(\beta\right)\label{eq:2.8}
\end{eqnarray}

Strategy for choosing grid points $z\left[1\right],z\left[2\right],........,z\left[m\right]$
depend on the nature of the data at hand�. 

For the measurement error model if we assume additive relationship
between precise and imprecise measurements as

\begin{equation}
W=X+u\;;E\left(u\right)=0,Var\left(u\right)=\tau^{2}\label{eq:2.9}
\end{equation}

Then 
\[
E\left(W|X\right)=X,\;\textrm{and }Var\left(W|X\right)=\tau^{2}
\]

When $\tau=0$, the distribution of $W|X,D$ is degenerated at the
point $x$. We simply take the grid-points to be the observed $x$
values.

When $\tau>0$,in the case of external validation there are competing
desiderata in choosing the $m$ grid-points. On the one hand, we want
m to be small in hopes of getting reasonable accuracy in estimating
$g$.� But on the other hand, to properly update the $x_{i}$ values,
we need a fine enough grid to correctly capture the shape of $f(w_{i}|x_{i},d_{i})$
as a function of $x_{i}$.�According to \cite{Gustafson_2002} we
will choose an equally spaced grid, with 
\[
z\left[1\right]=\min_{i}\left\{ w_{i}\right\} -2.5\tau,z\left[m\right]=\max_{i}\left\{ w_{i}\right\} +2.5\tau\textrm{ with a spacing of \ensuremath{\tau}/4.}
\]
 

\section{MCMC Scheme}

Our main interest is to draw samples from $f(\theta,\beta,\lambda|w,d)$
as presented in (\ref{eq:2.8}) using MCMC scheme.

For the measurement error model 
\[
W\mid X\sim N\left(X,\tau^{2}\right)
\]

$X\mid D,\beta,\lambda$ is approximated through $z\left[\theta\right]$ 

\[
logit\left\{ \mathbb{P}\left(D=1\mid X\right)\right\} =-3+0.5\exp\left(X\right)
\]

So, true value of $\beta$ is 0.5. The prior for $\beta$ is 
\[
\beta\sim N\left(0,100^{2}\right)
\]

Here it is not possible to specify a conditionally conjugate prior.
Also we assumed the candidate distribution to be normal which is symmetric.
So, we have used Metropolis sampling technique to replace draws from
the exact full conditional distribution with the draw from a candidate
distribution followed by an accept or reject state. 

\begin{algorithm}
\begin{algorithmic}[1]
\State {Initialize $\boldsymbol{\theta}^{\left(0\right)}=\left(\theta_{1}^{\left(0\right)},\theta_{2}^{\left(0\right)},.......,\theta_{p}^{\left(0\right)}\right)$}
\For{$s \gets 1$ to $S$} 
	\For{$j \gets 1$ to $p$}  
        \State {sample $\theta_{j}^{*}\sim q_{j}\left(\theta_{j}\mid\theta_{j}^{\left(s-1\right)}\right)$ for symmetric $q_j$}
		\State {set  $\boldsymbol{\theta}^{*}=\left(\theta_{1}^{\left(s\right)},...,\theta_{j-1}^{\left(s\right)},\theta_{j}^{*},\theta_{j+1}^{\left(s-1\right)},...,\theta_{p}^{\left(s-1\right)}\right)$ }
		\State{set $\displaystyle R=\frac{f\left(\boldsymbol{Y}\mid\theta^{*}\right)\pi\left(\theta^{*}\right)}{f\left(\boldsymbol{Y}\mid\theta^{\left(s-1\right)}\right)\pi\left(\theta^{\left(s-1\right)}\right)}$ }
		\State {sample $U \sim \textrm{uniform}(0,1)$}
		\If {$U<R$}
			\State {set $\theta_{j}^{\left(s\right)}=\theta_{j-1}^{*}$}	
		\Else
			\State {set $\theta_{j}^{\left(s\right)}=\theta_{j}^{\left(s-1\right)}$}
		\EndIf
	\EndFor
\EndFor
\end{algorithmic}

\caption{\textbf{Metropolis Algorithm}}
\end{algorithm}


\section{Computation}

Let us assume that the precise measurement $X\sim N\left(\mu,\tau^{2}\right)$.
By equation (\ref{eq:2.9}) 
\[
W\mid X\sim N\left(X,\tau^{2}\right)
\]

The model for $D\mid X$ is 

\begin{equation}
logit\left\{ \mathbb{P}\left(D=1\mid X\right)\right\} =\alpha+\beta\exp\left(X\right)\label{eq:4.1}
\end{equation}

We are interested in the posterior distribution of $X,\beta\mid W,D$.

\begin{eqnarray}
f(x,\beta|w,d) & \propto & f\left(w\mid x,d,\beta\right)f\left(x,d,\beta\right)\nonumber \\
 & = & f\left(w\mid x,d\right)f\left(x\mid d,\beta\right)f\left(d,\beta\right)\nonumber \\
 & = & f\left(w\mid x,d\right)f\left(x\mid d\right)f\left(d\mid\beta\right)f\left(\beta\right)\nonumber \\
 & = & \left(\prod_{i=1}^{n}f\left(w_{i}\mid x_{i},d_{i}\right)f\left(x_{i}\mid d_{i}\right)f\left(d_{i}\mid\beta\right)\right)f\left(\beta\right)\label{eq:4.2}
\end{eqnarray}

where 
\[
f\left(w_{i}\mid x_{i},d_{i}\right)=\frac{1}{\sqrt{2\pi}\tau}\exp\left(-\frac{\left(w_{i}-x_{i}\right)^{2}}{2\tau^{2}}\right)
\]

The expressions of $f\left(x_{i}\mid d_{i}\right)$ and $f\left(d_{i}\mid\beta\right)$
are derived as follows.

The $m$ grid-points $z\left[j\right]$'s are created as mentioned
in Section 0.2. To each grid-point we associate probability $\lambda_{j}$,
where $\left(\lambda_{1},...,\lambda_{m}\right)\sim Dirichlet\left(1,1,...,1\right)$.
We approximate the unobserved precise covariate $X$ by the grid-points.
So, $f\left(x\mid d\right)$ is a PMF with the $m$ grid-points as
the support.

\begin{equation}
f\left(x_{j}\mid d_{j}\right)=\frac{\lambda_{j}f\left(w_{j}\mid x_{j}\right)}{\sum_{j=1}^{m}\lambda_{j}f\left(w_{j}\mid x_{j}\right)},\;j=1\left(1\right)m\label{eq:4.3}
\end{equation}

Note that here we have $m$ $x_{j}$'s and $w_{j}$'s, whereas equation
(\ref{eq:4.2}) has $n$ many $x_{i}$'s and $w_{i}$'s. To overcome
this issue we simulate $w_{j}$ from $N\left(z\left[j\right],\tau^{2}\right),\;j=1,2,...,m$
and calculate $f\left(x_{j}\mid d_{j}\right)$. Then we generate a
sample of size $n$ with replacement from the distribution given by
(\ref{eq:4.3}).

The posterior of $X$ is given by 
\begin{equation}
f\left(x\mid w,d\right)\propto\prod_{i=1}^{n}f\left(w_{i}\mid x_{i},d_{i}\right)f\left(x_{i}\mid d_{i}\right)\label{eq:4.4}
\end{equation}

To draw samples from (\ref{eq:4.4}) we implemented MCMC using Metropolis
sampling algorithm.

From equation (\ref{eq:4.1})
\begin{eqnarray}
logit\left\{ \mathbb{P}\left(d_{i}=1\mid\bar{x_{i}}\right)\right\}  & = & \alpha+\beta\exp\left(\bar{x_{i}}\right)\nonumber \\
\implies\frac{\mathbb{P}\left(d_{i}=1\right)}{\mathbb{P}\left(d_{i}=0\right)} & = & \exp\left(\alpha+\beta\exp\left(\bar{x_{i}}\right)\right)\nonumber \\
\implies\frac{p_{i}}{1-p_{i}} & = & \exp\left(\alpha+\beta\exp\left(\bar{x_{i}}\right)\right)\;\textrm{where \ensuremath{p_{i}=}\ensuremath{\mathbb{P}\left(d_{i}=1\right)}}\nonumber \\
\implies p_{i} & = & \frac{\exp\left(\alpha+\beta\exp\left(\bar{x_{i}}\right)\right)}{1+\exp\left(\alpha+\beta\exp\left(\bar{x_{i}}\right)\right)}\label{eq:4.5}
\end{eqnarray}

where $\bar{x_{i}}$ is the post burn-in sample mean of $X_{i}$. 

So, 
\[
d_{i}\mid\beta\sim Ber\left(p_{i}\right)
\]

The posterior of $\beta$ is given by 
\begin{equation}
f\left(\beta\mid w,d\right)\propto\left(\prod_{i=1}^{n}f\left(d_{i}\mid\beta\right)\right)f\left(\beta\right)\label{eq:4.6}
\end{equation}

, where $f\left(\beta\right)$ is PDF of $N\left(0,100^{2}\right)$

To draw samples from (\ref{eq:4.6}), again we implemented Metropolis
sampling algorithm.

\begin{breakablealgorithm}
\caption{Computation}
\begin{algorithmic}[1]
\State{Set $\tau \in \left\{ 0,0.25,0.5,0.75\right\} $, $\mu=2$, $\alpha=-3$, $\beta=0.5$ and $n \in \left\{30,60,120,160\right\}$}
\State {Generate $\boldsymbol{X}=\left(X_{1},X_{2},..,X_{n}\right)$ such that $X_i \sim N \left(\mu,\tau^2 \right)$ ($n$ even)}
\State {Generate $\boldsymbol{D}=\left(D_{1},D_{2},..,D_{n}\right)$ with $n/2$ 0's and $n/2$ 1's such that}
	\If{$\displaystyle \frac{\exp\left(\alpha+\beta\exp\left(X\right)\right)}{1+\exp\left(\alpha+\beta\exp\left(X\right)\right)}<\frac{1}{2}$}
		\State {$D=0$}
	\Else
		\State {$D=1$}
	\EndIf
\State {Generate $\boldsymbol{W}=\left(W_{1},W_{2},..,W_{n}\right)$ such that}
\For{$i \gets 1$ to $n$}
	\State {$W_i \sim N \left(X_i,\tau^2 \right)$}
\EndFor
\State {Generate $m$ gridpoints $Z[1]<Z[2]<...<Z[m]$} such that
\For{$j \gets 1$ to $m$}
	\State {set $\displaystyle Z[j]=\min_{i}\left\{ w_{i}\right\} -2.5\tau+(j-1)\frac{\tau}{4}$}
\EndFor
\State {Draw $\left(\lambda_{1},\lambda_{2},...,\lambda_{m}\right) \sim Dirichlet (1, 1, ..., 1).$}
\State {Generate $\boldsymbol{W}=\left(W_{1},W_{2},..,W_{m}\right)$ and ${f}(\boldsymbol{w})=\left(f(W_1),f(W_2),..,f(W_m)\right)$such that}
\For{$j \gets 1$ to $m$}
	\State {Generate $W_j \sim N \left(Z[j],\tau^2 \right)$}
	\State {set $\displaystyle f(W_j)=\frac{1}{\tau\sqrt{2\pi}}\exp\left(-\frac{\left(W_j-Z\left[j\right]\right)^{2}}{2\tau^{2}}\right)$}
\EndFor
\State {Generate $\boldsymbol{p}=(p_1,p_2,...,p_m)$such that}
\For{$j \gets 1$ to $m$}
	\State {set $\displaystyle p_j=\frac{\lambda_{j}f\left(w_{j}\right)}{\sum_{k=1}^{m}\lambda_{j}f\left(w_{j}\right)}$}
\EndFor
\Function{CloseProb}{$a$}
	\State{return the probability of closest point of $a$}
\EndFunction
\Function{logPosteriorX}{$x$}
	\State{return log-Posterior of $x$}
\EndFunction
\State{set $S=10000$}
\State{Apply \emph{Metropolis Sampling} with initial value $\boldsymbol{X}^{(0)}=\boldsymbol{W}$, $p=n$ and candidate: $N \left(\boldsymbol{X}^{(s-1)},\tau^2 \right)$  }
\State{Consider samples $\boldsymbol{X}^{*}=\left(X_{1}^{*},X_{2}^{*},...,X_{n}^{*}\right)$ after burn-in period of 2000 samples}
\State{Obtain $\displaystyle \widehat{E\left(X\right)}=\frac{1}{8000n}\sum_{i=1}^{8000}\sum_{j=1}^{n}X_{ij}^{*}$, $\displaystyle \widehat{Var\left(X\right)}=\frac{1}{8000n}\sum_{i=1}^{8000}\sum_{j=1}^{n}\left(X_{ij}^{*}-\bar{X}_{j}^{*}\right)^{2}$, Mean distance $=\lVert \boldsymbol{X}-\bar{\boldsymbol{X}}^{*} \rVert$ where $\bar{\boldsymbol{X}}^{*} =\left(\bar{X}_{1}^{*},\bar{X}_{2}^{*},...,\bar{X}_{n}^{*}\right)$}
\Function{logPosteriorBeta}{$\beta$}
	\State{return log-Posterior of $\beta$}
\EndFunction
\State{Apply \emph{Metropolis Sampling} with initial value $\beta^{(0)}=50$, $p=1$ and candidate: $N \left(\beta^{(s-1)},100^2 \right)$  }
\State{Consider samples $\beta^{*}$ after burn-in period of 2000 samples}
\State{Obtain $\displaystyle \widehat{E\left(\beta\right)}=\frac{1}{8000}\sum_{i=1}^{8000}\beta_{i}^{*}$, $\displaystyle \widehat{MSE\left(\beta\right)}=\frac{1}{8000}\sum_{i=1}^{8000}\left(\beta_{i}^{*}-\beta\right)^{2}$, $\widehat{Var\left(\beta\right)}=\frac{1}{8000}\sum_{i=1}^{n}\left(\beta_{i}^{*}-\widehat{E\left(\beta\right)}\right)^{2}$}
\State{Compute $95\%$ credible interval  $\left( \displaystyle \widehat{E\left(\beta\right)}-z_{0.025}\sqrt{\frac{\widehat{Var\left(\beta\right)}}{8000}},\widehat{E\left(\beta\right)}+z_{0.025}\sqrt{\frac{\widehat{Var\left(\beta\right)}}{8000}}\right)$}
\State{Repeat (Step 2) to (Step 35) 10 times and obtain Coverage probability=Proportion of times credible interval contains the true value of $\beta$.} 
\end{algorithmic}
\end{breakablealgorithm}

\section{Results}

\subsection{Scenario A}
\begin{center}
\begin{table}[H]
\begin{centering}
\begin{tabular}{ccc}
\toprule 
$n_{i}$ & $\widehat{E\left(\beta\right)}$ & $\widehat{MSE\left(\beta\right)}$\tabularnewline
\midrule
\midrule 
15 & 0.5649 & 0.0128\tabularnewline
\midrule 
30 & 0.5594 & 0.0073\tabularnewline
\midrule 
60 & 0.5273 & 0.0034\tabularnewline
\midrule 
120 & 0.5116 & 0.0019\tabularnewline
\bottomrule
\end{tabular}
\par\end{centering}
\caption{Simulation results for scenario A. This scenario involves no measurement
error and various sample sizes (with $n_{1}=n_{2}$). The results
are based on 10 replicated data sets. The reported MSE is the MSE
for $E(\beta\mid data)$. }
\end{table}
\par\end{center}

\subsection{Scenario B}

Simulation results for scenario B for various values of $\tau$ are
considered. This scenario involves measurement error and fixed sample
size ( with $n_{1}=n_{2}=80$). The results are based on 10 replicated
data sets.
\begin{center}
\begin{table}[H]
\begin{centering}
\begin{tabular}{cccc}
\toprule 
$\tau$ & $\widehat{E\left(\beta\right)}$ & $\widehat{MSE\left(\beta\right)}$ & Coverage Probability\tabularnewline
\midrule
\midrule 
0.25 & 0.5256 & 0.0022 & 0.1\tabularnewline
\midrule 
0.5 & 0.4763 & 0.0052 & 0.3\tabularnewline
\midrule 
0.75 & 0.4490 & 0.0097 & 0.3\tabularnewline
\bottomrule
\end{tabular}
\par\end{centering}
\caption{Simulation results for scenario B for various values of $\tau$ are
considered, with $n_{1}=n_{2}=80$. Means and MSEs of $\hat{\beta}=E(\beta\mid data)$
are reported, along with the coverage probability (CP) of the 95\%
equal-tailed credible interval for $\beta$. The true value is $\beta$
= 0.5 throughout. }
\end{table}
\par\end{center}

\begin{center}
\begin{table}[H]
\begin{centering}
\begin{tabular}{cccc}
\toprule 
$\tau$ & $\widehat{E\left(X\right)}$ & $\sqrt{\widehat{Var\left(X\right)}}$ & Mean distance\tabularnewline
\midrule
\midrule 
0.25 & 1.9094 & 0.2543 & 2.3072\tabularnewline
\midrule 
0.5 & 1.9842 & 0.4835 & 8.0644\tabularnewline
\midrule 
0.75 & 1.9301 & 0.7559 & 6.4974\tabularnewline
\bottomrule
\end{tabular}
\par\end{centering}
\caption{Simulation results for scenario B for various values of $\tau$ are
considered, with $n_{1}=n_{2}=80$. Means and standard deviations
of $\hat{X}=E(X\mid data)$ are reported, along with the Mean distance
between the true. $X$ and predicted $X$. }
\end{table}
\par\end{center}

\section{Conclusion}
\begin{itemize}
\item In Scenario A involving no measurement error, we observe that the
estimated value of effect $\beta$ of the covariate is close to the
true value. The estimated Mean squared error of $\beta$ decreases
as the sample size increases and the posterior of $\beta$ yields
samples with better accuracy i.e., the effect of covariate can be
prominently observed as the sample size increases.
\item Scenario B involves measurement error for a fixed sample size (reasonably
large) the estimated Mean squared error of $\beta$ increases as the
variance of error increases and the expected value of $\beta$ is
not as close to the true value as in Scenario A. The method results
in a coverage probability of the covariate effect around 30\% when
applied on only 10 replicated datasets. It is expected that this probability
will increase as we increase the number of datasets. 
\item The mean distance between the true and estimated covariate values
is quite small and they approximately resemble the distribution of
the true covariate. Moreover, the estimated covariates from the posterior
can be used in cases where we need the covariates but do not actually
have it.
\end{itemize}

\section{Contribution}
\begin{itemize}
\item \textbf{\emph{Soumya Paul:}} Provided theoretical inputs,Writing code
(25\%), Preparing report(50\%), Presentation (35\%)
\item \textbf{\emph{Sankhadeep Mitra:}} Finding the paper (50\%), Writing
code (50\%), Presentation(35\%),Preparing report (25\%) 
\item \textbf{\emph{Sampriti Dutta:}} Finding the paper (50\%), Writing
code (25\%),Preparing report(25\%), Presentation (35\%) 
\end{itemize}

\bibliographystyle{ACM-Reference-Format}
\nocite{*}
\bibliography{Citation/BayesCite}

\end{document}
